\title{向量化排序算法的跨平台性能优化}
\author{马浩琨}
\date{Sep 16, 2026}
\maketitle
在数据库、搜索引擎、推荐系统以及图计算等场景中，排序操作始终扮演着底层核心的角色。即使在算法教科书中早已被视为「已解决」的问题，排序却在真实硬件约束下不断刷新性能边界。摩尔定律进入频率瓶颈后，芯片厂商把算力增益转向了并行维度，向量化执行成为必然选择。无论是 x86 上的 AVX-512、ARM 上的 SVE，还是 RISC-V 的可变长向量扩展，指令集碎片化与运行时差异对算法移植提出了新的挑战。本文的目标是给出一条可复制、可度量、可移植的向量化排序优化路线，让开发者能够在多平台上获得稳定、可预测的加速效果。\par
\chapter{背景与理论基础}
比较排序在理论上存在 $\Omega(n \log n)$ 的下界，而桶排序和基数排序则在整数或浮点键域上展示了线性时间 $O(n)$ 的潜力。向量化编程模型在不同架构上差异显著：x86 从 SSE 演进到 AVX-512，指令编码经历了 VEX 到 EVEX 的升级；ARM 从固定 128 位 NEON 进化为长度可变的 SVE/SVE2；RISC-V 的向量扩展 RVV 允许运行时决定向量寄存器长度；GPU 则以 SIMT 模型把线程束 warp 映射到大规模并行硬件。性能天花板往往并不来自计算单元，而是内存带宽、分支预测失败、TLB 抖动以及严格的对齐要求共同作用的结果。衡量指标不再局限于吞吐量 GB/s，还需关注 Cache 失效次数、能耗 J/GB 以及跨平台性能方差。\par
\chapter{跨平台向量化排序算法设计}
\section{基础并行基数排序 LSD}
最低有效位 LSD 基数排序通过「位」而非「值」进行分桶，避免了比较操作中的数据依赖。向量化改造的关键在于把顺序直方图累加替换为并行前缀和：先用向量指令一次性统计 16 或 32 个元素的桶计数，再通过并行扫描完成前缀和。随后的数据搬运可以使用「gather/scatter」指令，把分散在内存各处的元素一次性收集到连续的桶区域。整个过程不再依赖循环中的条件分支，而是用算术掩码与向量化比较完成决策，从而把控制流转化为数据流。\par
\section{Bitonic 与 Batchers 网络排序}
比较器网络天生适合 SIMD，因为每一步都是固定模式的 min/max 操作，没有数据相关的分支。实现时，可将 16 个 32 位元素装入一个 AVX2 寄存器，或把 8 个 64 位元素装入 NEON 寄存器；随后用「min/max」原语与「shuffle」指令在 lane 之间完成跨通道交换。Batchers 的奇偶归并网络把排序过程拆解为若干轮固定深度的比较-交换，每一轮都可映射到一条或数条向量指令，指令流水线得以充分填充。\par
\section{SIMT 版快速排序}
在 GPU 上，快速排序被重构为「块内排序 + 全局归并树」两阶段。块内排序利用共享内存执行小规模 Bitonic 网络；全局归并阶段则通过 CUDA 的 warp 同步原语把各块结果合并。关键在于避免 warp 内分支发散：所有线程在同一轮比较中执行相同的 min/max 操作，只有在归并树的不同层级才会出现控制流差异。CUDA Graph 可把多次内核启动固化为静态拓扑，进一步降低 CPU 端调度开销。\par
\section{混合策略}
真实负载中数组长度分布不均，需要运行时在不同算法间切换。小于等于 128 元素的数组直接使用 Bitonic 网络；中等规模（小于 $2^{20}$）采用 AVX2 并行基数排序；更大规模则启动多线程并行归并。运行时 CPU dispatch 通过 cpuid 或 getauxval 检测指令集支持，动态选择最优后端，避免 AVX-512 降频导致的负加速。\par
\chapter{跨平台工程实现}
\section{代码组织与编译策略}
项目采用 CMake 多后端结构：scalar.cpp 提供可移植参考实现；avx2.cpp 与 avx512.cpp 分别用 -mavx2、-mavx512f 等编译选项生成专用目标文件；neon.cpp 与 sve.cpp 针对 ARM 指令集；cuda.cu 负责 GPU 路径。运行时特征探测在程序启动时完成一次，把结果缓存到全局标志位，后续排序调用只需查表即可完成分发。\par
\section{数据布局与对齐}
强制 64 字节对齐能让向量加载落入单一 Cache Line，避免跨界拆分。C++17 的 alignas(64) 或 \_{}mm\_{}malloc 可满足这一需求。结构数组 AoS 改写为数组结构 SoA，把结构体成员拆成独立数组，可把 SIMD 访存效率提升一倍以上，因为相邻元素在内存中连续存放。\par
\section{内存带宽优化}
软件预取指令 \_{}mm\_{}prefetch 或 \_{}\_{}builtin\_{}prefetch 可在数据真正使用前若干周期发出访存请求，隐藏 DRAM 延迟。多路并行扫描则用 2 到 4 个游标交错读写同一数组，把访存流水线填满的同时降低 TLB 压力。实验显示，预取距离设置为 4 到 8 个 Cache Line 时，带宽利用率可从 60\%{} 提升至 85\%{}。\par
\section{分支消除与直方图优化}
查表法把条件判断替换成掩码运算：比较结果直接生成全 0 或全 1 的掩码向量，再与待选数据进行按位与运算。桶计数直方图若使用普通数组，会因多线程写冲突导致 Cache Line 乒乓；改用向量寄存器暂存局部直方图，最后再做一次跨寄存器归约，可把冲突降至最低。\par
\section{异构协同}
CPU 负责直方图统计与分桶索引生成，GPU 负责大规模元素搬运与最终归并。CUDA Graph 把多次小内核合并为单一图执行，消除了启动延迟。OpenMP 4.5/5.0 的 target 指令可把同一份代码同时 offload 到 GPU 与 DSP，无需维护两套代码库。\par
\chapter{性能评测与调优}
\section{测试环境与数据集}
实验覆盖 Intel ICL-8350、AMD Zen4、AWS Graviton3、Apple M2 以及 NVIDIA A100、AMD MI210。数据集包括均匀随机 32/64 位整数、Zipf 长尾分布，以及真实日志与点击流。指标不仅记录吞吐量，还统计带宽利用率、Speed-up 相对于 STL std::sort 的倍数，以及跨平台性能变异系数 CV。\par
\section{典型调优案例}
在 Intel 平台上，盲目使用 AVX-512 导致核心频率下降，反而慢于 AVX2 多线程版本；降级策略改回 AVX2 后，实测带宽利用率回升 18\%{}。在 ARM 平台上，NEON 的 gather/scatter 吞吐远低于标量循环；改用「转置 + 向量化比较」的混合方案后，性能反超原始标量实现 2.3 倍。\par
\chapter{工程落地 checklist}
统一 SIMD 抽象层可选用 xsimd、highway 或 simde，避免重复编写多套 intrinsics。持续集成需覆盖 QEMU 用户态、AWS ARM 实例与 NVIDIA Docker 镜像，确保每次提交都在真实硬件上回归。性能监控采用 Google Benchmark 搭配 benchstat 做统计显著性检验，避免偶然波动误判。文档中需明确列出各后端的指令集需求、编译 flag 以及运行时 dispatch 逻辑，降低新人上手成本。\par
\chapter{未来展望}
Intel AMX、ARM SME 与 RISC-V RVV 1.0 正在引入矩阵运算原语，排序算法有望直接利用矩阵乘法单元完成并行归并。端到端编译器 MLIR/LLVM 已出现实验性向量化排序 Pass，可自动生成多后端代码。未来近存计算 PIM 与 CXL 内存池将把排序算子下推到内存控制器，彻底改变「先搬数据再计算」的传统范式。\par
\chapter{结论}
向量化排序的极致性能来自「算法选择、数据布局、指令集感知」三位一体的协同设计。开发者不应盲目追求理论峰值，而应在真实负载上持续 profiling、度量、调优，才能把硬件潜能转化为业务收益。\par
\chapter{附录}
\section{A. Bitonic 16 路 AVX2 关键代码片段}
\begin{lstlisting}[language=c]
// 加载 16 个 uint32_t 到 ymm0
__m256i v = _mm256_load_si256(reinterpret_cast<const __m256i*>(ptr));
// 与交换距离为 8 的 lane 做比较
__m256i vswap = _mm256_permute4x64_epi64(v, _MM_SHUFFLE(1,0,3,2));
__m256i min1 = _mm256_min_epu32(v, vswap);
__m256i max1 = _mm256_max_epu32(v, vswap);
// 再做一次 shuffle 完成跨 128 位边界交换
__m256i min2 = _mm256_shuffle_epi32(min1, _MM_SHUFFLE(2,3,0,1));
__m256i max2 = _mm256_shuffle_epi32(max1, _MM_SHUFFLE(2,3,0,1));
\end{lstlisting}
这段代码首先使用 \_{}mm256\_{}load\_{}si256 把 16 个 32 位无符号整数加载到 256 位 AVX2 寄存器。随后的 \_{}mm256\_{}permute4x64\_{}epi64 指令把高 128 位与低 128 位做整体交换，实现距离为 8 个元素的跨 lane 比较。min/max 原语完成无分支的比较-交换；第二次 shuffle 则在 128 位边界内部再做一次 2 元素交换，把排序网络向前推进一层。整个过程没有任何标量循环，全部映射为 4 条向量指令，延迟固定且可流水。\par
\section{B. 性能数据表格（原始 CSV 节选）}
\begin{lstlisting}
platform,algorithm,n_elements,throughput_GB,speedup_vs_std
Intel_ICL,avx2_radix,1e7,18.4,4.2
Graviton3,neon_bitonic,1e7,9.7,2.8
A100,cuda_merge,1e8,92.1,12.5
\end{lstlisting}
\section{C. 参考文献}
Intel® 64 and IA-32 Architectures Optimization Reference Manual\par
Arm® Architecture Reference Manual Supplement SVE\par
CUDA C Programming Guide v12.0\par
Highway: Performance-portable SIMD library\par
xsimd: Modern, portable SIMD intrinsics wrapper\par
